% ============================================================
%  Introducción y objeto
% ============================================================
\section{Introducción}
\label{sec:introduccion}

En la industria de la fabricación de dispositivos eléctricos, la seguridad
eléctrica es un aspecto crítico para garantizar la protección de los usuarios.
La norma IEC~60335 establece los ensayos de seguridad que deben cumplir los
aparatos electrodomésticos y similares antes de su comercialización
\cite{iec60335}. Los métodos tradicionales de prueba, que dependen en gran medida
de la intervención manual del operario, pueden dar lugar a variaciones en los
resultados, prolongar los tiempos de prueba y aumentar los riesgos tanto para el
personal como para los dispositivos bajo evaluación.

El presente informe describe el diseño, la construcción y la verificación de un
equipo portátil y semiautomatizado para ejecutar esos ensayos en final de línea,
los resultados obtenidos en el banco de pruebas y las conclusiones del trabajo.

\section{Objeto}
\label{sec:objeto}

La seguridad eléctrica se ha convertido en un pilar indispensable para la
confiabilidad y comercialización de los equipos de uso doméstico e industrial. En
este contexto, surge la necesidad de innovar en los procesos de verificación,
superando las limitaciones de los métodos tradicionales que dependen en gran
medida de la intervención humana. Este proyecto tiene por objeto el diseño y la
construcción de un dispositivo semiautomatizado capaz de ejecutar, con precisión
y repetibilidad, las principales pruebas de seguridad definidas en la norma
internacional IEC~60335 ---rigidez dieléctrica, continuidad de tierra, medición
de corrientes de fuga y consumo---, garantizando así un control de calidad más
eficiente y seguro.

\vspace{0.3cm}\noindent El desarrollo se complementa con una interfaz de software
que integra los distintos subsistemas, procesa los datos obtenidos y facilita la
interacción con el operario. De este modo, se busca no solo reducir los tiempos de
prueba y minimizar los riesgos, sino también aportar una herramienta que eleve los
estándares de confiabilidad y trazabilidad en el control de calidad de productos
eléctricos.
